% T13 -- Which artefact produces which headline number.
\begin{minipage}{\textwidth}
\centering\footnotesize
\captionof{table}{The origin of each number quoted in the report}
\label{tab:provenance}
\begin{tabular}{@{}>{\raggedright\arraybackslash}p{4.5cm}ll>{\raggedright\arraybackslash}p{6.0cm}@{}}
\toprule
\textbf{Number quoted in the report} & \textbf{Data} & \textbf{Iter.} & \textbf{Which run produced it} \\
\midrule
16.87\,\% first signal              & \sraw & 1  & single-notebook Mistral-7B QLoRA run, 100 steps \\
48.7\,\% $\rightarrow$ 64.4\,\%     & \sraw & 2  & same notebook, keeping the best saved copy \\
$r=-0.971$ on the switch rate       & \textsc{raw day}, 500 people & 2 & per-person analysis of the 1{,}173-person grading \\
$+1.1$ at the 80/20 cut             & \sfilt & 4  & training run on the filtered day \\
No cut beats the rule               & \smob & 5  & one trained model, seven grading passes \\
$-3.0$ from shuffling the hours     & \smob & 6  & re-grading the same trained copy \\
$+2.6$ at 18--20h                   & \smob & 6  & the targeted-sampling run, whole day graded \\
$27.4$ separation of the derived hours & \smob & 7 & analysis of the kinds of user on the training file \\
$+3.7$, first run above the rule    & \smob & 8  & routing each kind of user to its own numbers \\
$+4.2$ at 1{,}900 training people   & \stwok & 9  & scaling run on the 2{,}000-person day \\
$-1.6$ on stretched days            & \slev & 9  & the 200-record levelling run \\
$+4.3$, flat from 5{,}000 to 21{,}000 & \seleven & 10 & four-size scaling notebook, constant training time \\
The resampling intervals            & \seleven & 10 & 2{,}000 paired resamples over the test people \\
0 of 7 weekday trainings finished   & \seleven & 11 & seven-day notebook; the session crashed after 5\,h \\
\bottomrule
\end{tabular}

\figtext{%
\figpara{Table~\ref{tab:provenance}: data provenance mapping for experimental results.}{Every number quoted anywhere in the body can be
traced back to one run through Table~\ref{tab:provenance}. The \textbf{Data} column names the set
of people it was measured on, described in \textbf{Table~\ref{tab:datasets}}, and
\textbf{Iter.} is the numbered iteration of \textbf{Appendix~\ref{app:timeline}}.}

\figpara{Table~\ref{tab:provenance}: reproducibility tracking across notebook runs.}{Two rows in Table~\ref{tab:provenance} measured on
different sets of people are not comparable on their absolute scores, only on
their distances to the rule of that row's own data. Reading Table~\ref{tab:provenance} alongside
\textbf{Table~\ref{tab:datasets}} is what makes that check possible.}}

\end{minipage}
